\section*{Bab \ref{ch:g4:fractions} --- Pecahan Pertama}

\begin{solution}{exo:g4:fractions:1}
Pizza: tersisa $\frac56$. Cokelat: $\frac38$ dimakan. Bendera:
$\frac13$ diberi warna.
\end{solution}

\begin{solution}{exo:g4:fractions:2}
$\frac35$ mewarnai tiga bagian, $\frac25$ hanya dua: $\frac35$
adalah warna yang lebih besar.
\end{solution}

\begin{solution}{exo:g4:fractions:3}
Setengah dari $18$: $9$. \omterm{def:g3:division:half}{Seperempat} dari $20$: $5$. Sepertiga dari
$21$: $7$. $\frac34$ dari $20$: tiga \omterm{def:g3:division:half}{seperempat}, $3 \times 5 = 15$.
\end{solution}

\begin{solution}{exo:g4:fractions:4}
$\frac13$: satu langkah; $\frac33$: tepat di $1$; $\frac53$: dua
langkah melewati $1$. \omterm{def:g4:fractions:def}{Pecahan} yang mendarat tepat di $2$ adalah $\frac63$.
\end{solution}

\begin{solution}{exo:g4:fractions:5}
$\frac38 < \frac68$; \quad $\frac55 = 1$; \quad $\frac74 > 1$;
\quad $\frac12 = \frac24$.
\end{solution}

\begin{solution}{exo:g4:fractions:6}
$\frac26 + \frac36 = \frac56$; \quad
$\frac58 - \frac28 = \frac38$; \quad
$\frac14 + \frac24 = \frac34$; \quad
$1 = \frac33$, jadi $1 - \frac13 = \frac23$.
\end{solution}

\begin{solution}{exo:g4:fractions:7}
$\frac54 = 1 + \frac14$; \quad
$\frac73 = 2 + \frac13$; \quad
$\frac92 = 4 + \frac12$.
\end{solution}

\begin{solution}{exo:g4:fractions:8}
Setengah jam: $30$ menit. \omterm{def:g3:division:half}{Seperempat}: $15$ menit. Tiga perempat:
$45$ menit. Sepertiga: $20$ menit.
\end{solution}

\begin{solution}{exo:g4:fractions:9}
$\frac14$ dari $24$ adalah $6$ murid yang memakai kacamata;
$24 - 6 = 18$ murid tidak memakai.
\end{solution}

\begin{solution}{exo:g4:fractions:10}
Bersama: $\frac38 + \frac48 = \frac78$. Tersisa: $\frac18$. Sandi
makan lebih banyak ($4$ perdelapan melawan $3$).
\end{solution}

\begin{solution}{exo:g4:fractions:11}
Memang mungkin: setengah pizza kecil bisa lebih sedikit daripada
\omterm{def:g3:division:half}{seperempat} pizza raksasa. \omterm{def:g4:fractions:def}{Pecahan} membandingkan bagian \emph{dari
keseluruhan yang sama}; sebelum membandingkan $\frac12$ dan
$\frac14$ ``dari sesuatu'', kita harus tahu bahwa kedua sesuatu itu sama.

\end{solution}
